%------------------------------------------------------------------------------
% 3.5 نمذجة منظومة الوكلاء والعلاقات التفاعلية بينهم
% NEW section (2026-10) — the direct response to the committee's objection:
% agents are modelled as INTERNAL system components with an explicit
% interaction protocol (spec matrix, shared AgentState, orchestration
% patterns, and a full UML sequence diagram fulfilling Gantt task 11.4.1).
%------------------------------------------------------------------------------
\section{نمذجة منظومة الوكلاء والعلاقات التفاعلية بينهم}
\label{sec:agent-modeling}

بعد أن حدّد القسم السابق حالات الاستخدام من منظور الفاعلين البشريين الخارجيين، يتناول هذا القسم الوجه الآخر للنمذجة: البنية الداخلية لمنظومة الوكلاء التي تنفذ تلك الحالات داخل حدود النظام. فالوكلاء الخمسة --- المشرف، وقواعد البيانات، والاسترجاع، والتحليل، والتقارير --- ليست فاعلين خارجيين، بل هي مكونات برمجية متخصصة تنسق فيما بينها عبر إطار \EN{LangGraph} \cite{langgraph2024} وفق بروتوكول تفاعل صريح قابل للتحقق، وهو ما تمليه أدبيات الأنظمة متعددة الوكلاء القائمة على النماذج اللغوية الكبيرة \cite{wu2023autogen}. ويستند هذا القسم إلى أربع ركائز متكاملة: التوصيف الهيكلي لكل وكيل، ثم نموذج الحالة المشتركة الذي يحكم التخاطب بينهم، ثم أنماط التنسيق الأربعة التي تحكم مسارات التنفيذ، وأخيراً نمذجة التسلسل التفاعلي الشامل لاستعلام مؤسسي متكامل.

\subsection{التوصيف الهيكلي للوكلاء}
\label{subsec:agent-specification}

يُعرَّف كل وكيل في نظام «بيان» وفق منهجية التوصيف الهيكلي المعتمدة في أدبيات الوكلاء الأذكياء، والتي تحدد هوية الوكيل من أربع زوايا: دوره التنظيمي داخل المنظومة، ومدخلاته السياقية التي يدركها، وأدواته وصلاحياته المتاحة له، ومخرجاته والوجهة التي يسلمها إليها. ويضمن هذا التوصيف الدقيق أن حدود المسؤولية بين الوكلاء واضحة لا تتقاطع، وأن لكل وكيل غرضاً واحداً رئيساً يُتقنه، بما يحقق مبدأ التخصص والتقارب في تصميم الأنظمة متعددة الوكلاء. ويعرض الجدول \ref{tab:agent-spec} مصفوفة التوصيف الكاملة للوكلاء الخمسة.

\begin{table}[H]
\centering
\caption{مصفوفة التوصيف الهيكلي لوكلاء نظام «بيان» الخمسة.}
\label{tab:agent-spec}
\scriptsize
\renewcommand{\arraystretch}{1.35}
\begin{tabularx}{\textwidth}{|>{\raggedright\arraybackslash}p{2.0cm}|>{\raggedright\arraybackslash}p{2.5cm}|>{\raggedright\arraybackslash}X|>{\raggedright\arraybackslash}X|>{\raggedright\arraybackslash}X|}
\hline
\textbf{الوكيل} & \textbf{الدور التنظيمي} & \textbf{المدخلات السياقية} & \textbf{الأدوات والصلاحيات} & \textbf{المخرجات والوجهة التالية} \\
\hline
الوكيل المشرف \EN{(Supervisor)}
& منسق مركزي وموجّه النوايا \EN{(Orchestrator \& Router)}
& سؤال المستخدم بالعربية، وسجل الجلسة، وصلاحيات المستخدم \EN{(RBAC)}
& مصنّف النوايا \EN{(Intent Classifier)}، وقارئ حالة الجلسة \EN{(AgentState)}، ومحرك الحواف المشروطة
& خطة التوجيه \EN{(Routing Plan)}: إسناد المهمة لوكيل واحد أو سلسلة وكلاء متتابعة، ثم استلام المخرج النهائي وتسليمه لواجهة المحادثة. \\
\hline
وكيل قواعد البيانات \EN{(SQL)}
& مترجم ومنفّذ الاستعلامات المهيكلة
& السؤال المهيكل، ومخطط \EN{Smart Views} المسموح لدور المستخدم
& أداة ربط المخطط \EN{(Schema Linker)}، ومولّد \EN{SQL}، ومنفّذ القراءة المعزول \EN{(Read-Only Executor)}، ووحدة التصحيح الذاتي
& جدول بيانات خام \EN{(JSON/DataFrame)} مع الاستعلام المنفَّذ؛ يُسلَّم للوكيل المشرف أو لوكيل التحليل. \\
\hline
وكيل الاسترجاع \EN{(RAG)}
& باحث اللوائح والمستندات غير المهيكلة
& السؤال النصي، وفلاتر الصلاحيات الوثائقية
& نموذج التضمين \EN{multilingual-e5-large}، ومحرك البحث \EN{Qdrant} \cite{qdrant2024}، وأداة استخراج التوثيق \EN{(Citation Extractor)}
& مقاطع موثقة بأرقام الصفحات والمصادر؛ تُسلَّم للمشرف أو تُمرَّر كمعامل شرطي لوكيل \EN{SQL}. \\
\hline
وكيل التحليل \EN{(Analyst)}
& محلل إحصائي ومكتشف الأنماط والشذوذ
& البيانات الخام المستخرجة من وكيل \EN{SQL}، وهدف التحليل
& مكتبات التحليل الإحصائي، وكاشف القيم الشاذة \EN{(Outlier Detector)}، ومولّد توصيف الرسوم البيانية
& مؤشرات إحصائية وتنبيهات شذوذ وبيانات رسوم؛ تُسلَّم لوكيل التقارير أو المشرف. \\
\hline
وكيل التقارير \EN{(Reporter)}
& منسّق المخرجات ومولّد الوثائق
& مخرجات وكلاء \EN{SQL} و\EN{RAG} والتحليل من الحالة المشتركة
& قوالب التنسيق العربية \EN{RTL}، ومحرك التصدير \EN{(PDF/Excel/CSV)}، ومنسّق الجداول العربية
& تقرير نهائي منسق قابل للعرض والتنزيل؛ يُعرض للمستخدم عبر الوكيل المشرف. \\
\hline
\end{tabularx}
\end{table}

ويلاحظ من المصفوفة أن الوكيل المشرف يحتكر مسؤولية التوجيه والتحقق من الصلاحيات، بينما يحتفظ كل وكيل متخصص باستقلالية تنفيذية داخل نطاق أدواته فقط؛ فوكيل \EN{SQL} مثلاً لا يملك صلاحية الكتابة في قاعدة البيانات إطلاقاً، إذ ينفذ استعلاماته عبر منفّذ القراءة المعزول حصراً، وهو قرار تصميمي يعزز الأمان المؤسسي ويجعل كل عملية قابلة للإسناد والمساءلة عبر سجل التدقيق.

\subsection{نموذج الحالة المشتركة وبروتوكول التخاطب}
\label{subsec:agent-state-schema}

لا يتبادل وكلاء نظام «بيان» رسائل ثنائية عشوائية فيما بينهم؛ بل يعتمدون في تنسيقهم على نمط تمرير الحالة \EN{(State-Passing)} الذي يوفره إطار \EN{LangGraph}، حيث يمرَّر كائن حالة موحد \EN{(AgentState)} عبر عقد المخطط، فتقرأ كل عقدة الحقول التي تحتاجها من الحالة، وتحدّث الحقول التي تنتجها، ثم تمرر الحالة المحدثة إلى العقدة التالية عبر حافة مشروطة \EN{(Conditional Edge)} يحددها الوكيل المشرف بناءً على خطة التوجيه. وتضمن هذه البنية ثلاث خصائص جوهرية: أولاً، اتساق المعلومات إذ لا يمكن لوكيلين العمل على نسختين متضاربتين من البيانات؛ وثانياً، قابلية التتبع الكامل لأن كل تحديث للحالة يُسجَّل في مسار التنفيذ \EN{(Agent Trace)}؛ وثالثاً، قابلية الاستئناف عند الفشل لأن الحالة تحمل كل ما يلزم لإعادة المحاولة دون فقدان السياق. ويلخص الجدول \ref{tab:agent-state-schema} الحقول الأساسية في كائن الحالة المشتركة مع الوكيل القارئ والوكيل الكاتب لكل حقل.

\begin{table}[H]
\centering
\caption{الحقول الأساسية في كائن الحالة المشتركة \EN{(AgentState)} لوكلاء نظام «بيان».}
\label{tab:agent-state-schema}
\scriptsize
\renewcommand{\arraystretch}{1.35}
\begin{tabularx}{\textwidth}{|>{\raggedright\arraybackslash}p{3.9cm}|>{\raggedright\arraybackslash}X|>{\raggedright\arraybackslash}p{3.4cm}|>{\raggedright\arraybackslash}p{3.2cm}|}
\hline
\textbf{الحقل} & \textbf{الدلالة والوظيفة} & \textbf{الوكيل القارئ} & \textbf{الوكيل الكاتب} \\
\hline
\EN{session\_id} & المعرّف الفريد للجلسة الجارية & جميع الوكلاء & الواجهة الخلفية \EN{(FastAPI)} \\
\hline
\EN{user\_query} & نص سؤال المستخدم بالعربية كما ورد & المشرف، ووكيل \EN{SQL}، ووكيل \EN{RAG} & الواجهة الخلفية \\
\hline
\EN{user\_role\_rbac} & دور المستخدم ومستوى صلاحيته الأمنية & المشرف، ووكيل \EN{SQL}، ووكيل \EN{RAG} & وحدة \EN{RBAC} \\
\hline
\EN{routing\_plan} & خطة التوجيه: ترتيب الوكلاء ومسار التنفيذ & جميع الوكلاء & الوكيل المشرف \\
\hline
\EN{target\_smart\_views} & قائمة الرؤى الذكية المسموح بها لدور المستخدم & وكيل \EN{SQL} & الوكيل المشرف \\
\hline
\EN{generated\_sql} & الاستعلام المولَّد والمنفَّذ مع نتيجته & وكيل التحليل، وسجل التدقيق & وكيل \EN{SQL} \\
\hline
\EN{retry\_count} & عدد محاولات التصحيح الذاتي المتحققة & وكيل \EN{SQL} & وكيل \EN{SQL} \\
\hline
\EN{raw\_dataset} & البيانات الخام المستخرجة من قاعدة البيانات & وكيل التحليل، ووكيل التقارير & وكيل \EN{SQL} \\
\hline
\EN{retrieved\_chunks\_citations} & المقاطع المسترجعة وتوثيق مصادرها & المشرف، ووكيل التقارير & وكيل \EN{RAG} \\
\hline
\EN{statistical\_insights\_outliers} & المؤشرات الإحصائية والقيم الشاذة المكتشفة & وكيل التقارير، والمشرف & وكيل التحليل \\
\hline
\EN{formatted\_report\_artifact} & التقرير النهائي المنسق القابل للتصدير & الوكيل المشرف (للتسليم) & وكيل التقارير \\
\hline
\EN{agent\_trace\_steps} & مسار تنفيذ الوكلاء خطوة بخطوة & وحدة التدقيق، وواجهة العرض & جميع الوكلاء (إلحاقياً) \\
\hline
\end{tabularx}
\end{table}

\subsection{أنماط العلاقات وتنسيق الوكلاء}
\label{subsec:agent-orchestration-patterns}

تتبع العلاقات بين الوكلاء في نظام «بيان» نمطاً عاماً واحداً هو \textbf{التنسيق الهجين الموجه بالحالة \EN{(Stateful Hybrid Supervisor-Pipeline Pattern)}}: تنسيق مركزي عبر الوكيل المشرف، تتشكل داخله سلاسل تنفيذ متتابعة أو متوازية بحسب طبيعة الاستعلام، وتُدار جميعها من خلال كائن الحالة المشترك. ويتفرع من هذا النمط العام أربعة أنماط تفاعل رئيسة تحكم مسارات التنفيذ:

\textbf{أولاً: نمط التفويض الأحادي المباشر \EN{(Single-Agent Direct Routing)}:} وهو المسار الافتراضي للاستعلامات البسيطة، إذ يحلل الوكيل المشرف نية المستخدم ويفحص صلاحيته ثم يوجه الطلب إلى وكيل واحد فقط: وكيل \EN{SQL} للاستعلامات المهيكلة المباشرة (مثل: «كم عدد الطلاب المسجلين هذا الفصل؟»)، أو وكيل \EN{RAG} للأسئلة الوثائقية المباشرة (مثل: «ما شروط منح الإجازة السنوية؟»). وبعد إعادة الوكيل المتخصص نتيجته إلى المشرف عبر الحالة المشتركة، يسلّمها المشرف إلى واجهة المحادثة مباشرة دون المرور ببقية الوكلاء، بما يقلل زمن الاستجابة ويحافظ على بساطة المسار.

\textbf{ثانياً: نمط السلسلة التحليلية المتتابعة \EN{(Sequential Analytical Pipeline)}:} ويُستدعى عند طلب تقارير تحليلية مخصصة (حالة الاستخدام \EN{UC-03})، إذ يبني المشرف سلسلة تنفيذ متتابعة: يستدعي وكيل \EN{SQL} أولاً لاستخراج البيانات الخام عبر طبقة الرؤى الذكية، ثم يمررها إلى وكيل التحليل الذي يحسب المؤشرات الإحصائية ويكتشف القيم الشاذة والاتجاهات، ثم يسلم مخرجاته إلى وكيل التقارير الذي ينسق التقرير النهائي بصيغ قابلة للتصدير. وتضمن تتابعية السلسلة أن كل وكيل يعمل على مدخلات جاهزة ومعتمدة من سابقه، مع إمكانية تدقيق كل خطوة على حدة في مسار التنفيذ.

\textbf{ثالثاً: نمط التكامل الهجين بين اللوائح والبيانات \EN{(Hybrid Policy-to-Data Synthesis)}:} وهو أكثر الأنماط تعبيراً عن القيمة المضافة للمنظومة، ويُستدعى للاستفسارات المركبة التي تدمج المعيار التنظيمي بالبيانات التشغيلية، مثل: «أظهر الطلاب المتجاوزين لنسبة الغياب المسموحة وفق لائحة الجامعة». إذ يوجّه المشرف طلباً فرعياً أول إلى وكيل \EN{RAG} لاستخراج النسبة المسموحة من اللائحة (مثلاً 25\%) مع توثيق مصدرها، ثم تمرَّر هذه القيمة عبر الحالة المشتركة كمعامل شرطي \EN{(Parameter)} إلى وكيل \EN{SQL} الذي يبني استعلامه على أساسها (الطلاب الذين تجاوز غيابهم 25\%)، ثم تُكمل السلسلة عبر وكيل التحليل ووكيل التقارير إن لزم. وبهذا يجيب النظام عن سؤال واحد يجمع بين مصدرين مختلفين في الطبيعة، مع توثيق كامل للمصدر التشريعي للمقارنة.

\textbf{رابعاً: حلقة التصحيح الذاتي المشروطة \EN{(Conditional Self-Correction Loop)}:} وهي حلقة تكرارية داخلية ضمن وكيل \EN{SQL} تُنشَّط عند فشل الاستعلام المولَّد نحوياً أو رفضه أمنياً، إذ يغذي الوكيل رسالة الخطأ راجعةً إلى النموذج اللغوي مقرونة بأمثلة تصحيحية \EN{Few-Shot} لإعادة توليد الاستعلام، بحد أقصى ثلاث محاولات مسجلة في حقل \EN{retry\_count}. فإذا استنفدت المحاولات دون نجاح، أعاد الوكيل للمشرف اعتذاراً منضبطاً يوضح سبب التعذر ويقترح إعادة صياغة السؤال، بدلاً من عرض إجابة خاطئة أو غير مدعومة.

ويجسد الشكل \ref{fig:agent-state-graph} هذه الأنماط مجتمعة في مخطط آلة الحالة وتدفق العقد المشروطة عبر \EN{LangGraph}: إذ تظهر العقد الخمس، وكائن الحالة المشتركة بناقله الذي يقرأه ويحدّثه كل وكيل، والحواف المشروطة الصادرة عن المشرف بعناوينها، وحلقة التصحيح الذاتي، والمسار الهجين لتمرير المعامل، فضلاً عن الارتباط بمصادر البيانات وسجل التدقيق.

\begin{figure}[H]
  \centering
  \resizebox{\textwidth}{!}{%==============================================================================
% figures/tikz/agent-state-graph.tex
% Multi-Agent State & Orchestration Graph — new for Section 3.5 (2026-10).
%
% Shows (per the thesis agent-modeling section):
%   - The five agent nodes orchestrated by the Supervisor via conditional
%     edges (labels: استعلام بيانات / سؤال لوائح / تحليل وتقرير).
%   - The shared AgentState object + a dashed state bus: every node reads
%     and updates the common state (LangGraph state-passing model).
%   - The conditional self-correction loop inside the SQL agent (max 3
%     retries), the hybrid RAG->SQL parameter path, the analyst conditional
%     call, and the final return path to the supervisor.
%   - Data layers: PostgreSQL via Smart Views + Qdrant; Audit Log & Trace.
% All Arabic labels are wrapped in \textarabic{} (proven RTL fix inside
% TikZ LTR nodes). Depends on main.tex colors + \EN macro.
%==============================================================================
\resizebox{\textwidth}{!}{%
\begin{tikzpicture}[
  >=stealth,
  font=\small,
  pill/.style={
    rounded corners=8pt, draw=kauaccent, fill=kauaccent!8, text=kaudark,
    font=\small\bfseries, align=center, inner sep=6pt,
    minimum width=4.6cm, minimum height=0.62cm
  },
  supbox/.style={
    rectangle, rounded corners=5pt, draw=kauprimary, fill=kauprimary!12,
    text=kaudark, font=\small\bfseries, align=center,
    minimum width=5.6cm, minimum height=1.2cm, line width=1.2pt
  },
  agentbox/.style={
    rectangle, rounded corners=4pt, draw=kausecondary, fill=kausecondary!10,
    text=kaudark, font=\footnotesize\bfseries, align=center,
    minimum width=3.4cm, minimum height=1.2cm, inner sep=4pt
  },
  statebox/.style={
    rectangle, rounded corners=5pt, draw=kauprimary, fill=kauprimary!8,
    text=kaudark, font=\footnotesize\bfseries, align=center,
    minimum width=3.5cm, minimum height=1.3cm, dashed
  },
  datastore/.style={
    cylinder, shape border rotate=90, aspect=0.3,
    draw=kaudark, fill=kaudark!8, text=kaudark,
    font=\scriptsize, align=center, inner sep=2pt
  },
  auditbox/.style={
    rectangle, rounded corners=4pt, draw=kaudark!70, fill=kaudark!5,
    text=kaudark, font=\scriptsize\bfseries, align=center,
    minimum width=4.0cm, minimum height=0.9cm
  },
  outpill/.style={
    rounded corners=7pt, draw=kauaccent, fill=kauaccent!12, text=kaudark,
    font=\footnotesize\bfseries, align=center,
    minimum width=3.8cm, minimum height=0.75cm
  },
  arrowS/.style={->, >=stealth, line width=1.0pt, draw=kauprimary},
  arrowA/.style={->, >=stealth, line width=0.7pt, draw=kausecondary!90!black},
  arrowR/.style={->, >=stealth, line width=0.7pt, draw=kauaccent!90!black, dashed},
  arrowL/.style={->, >=stealth, line width=0.6pt, draw=kaudark!55, dotted},
  statbus/.style={draw=kauprimary!60, line width=0.75pt, dashed},
  lbl/.style={font=\tiny\bfseries, text=kaudark, fill=white, inner sep=1.2pt},
  lbls/.style={font=\tiny, text=kaudark!90, fill=white, inner sep=1.2pt},
]

  %==========================================================================
  % ENTRY: user query (top-centre) -> Supervisor
  %==========================================================================
  \node[pill] (query) at (2.2,1.0)
    {\textarabic{استعلام المستخدم بالعربية}};
  \node[supbox] (sup) at (2.2,-0.5)
    {\textarabic{الوكيل المشرف \EN{(Supervisor)}}\\[1pt]
     {\scriptsize\textarabic{تحليل النية + فحص \EN{RBAC} + التوجيه المشروط}}};
  \draw[arrowS] (query.south) -- (sup.north);

  %==========================================================================
  % AGENTS: SQL (right = first read) and RAG (left)
  %==========================================================================
  \node[agentbox] (sql) at (6.4,-3.3)
    {\textarabic{وكيل قواعد البيانات \EN{(SQL)}}\\[1pt]
     {\scriptsize\textarabic{عبر طبقة \EN{Smart Views}}}};
  \node[agentbox] (rag) at (0.6,-3.3)
    {\textarabic{وكيل الاسترجاع \EN{(RAG)}}\\[1pt]
     {\scriptsize\textarabic{بحث دلالي في \EN{Qdrant}}}};

  % Supervisor conditional dispatch edges
  \draw[arrowA] (sup.south east) -- (sql.north)
    node[lbls, pos=0.55, right=1pt] {\textarabic{استعلام بيانات}};
  \draw[arrowA] (sup.south west) -- (rag.north)
    node[lbls, pos=0.55, left=1pt] {\textarabic{سؤال لوائح}};

  % RAG -> Supervisor return (documented chunks)
  \draw[arrowR] (rag.north west)
    .. controls (-1.6,-2.4) and (-1.6,-1.4) .. (sup.south west)
    node[lbls, pos=0.30, left=2pt] {\textarabic{مقاطع موثقة}};

  % Hybrid pattern: RAG -> SQL parameter pass (regulatory threshold)
  \draw[arrowR] (rag.east) -- (sql.west)
    node[lbls, midway, above=2pt] {\textarabic{معامل من اللائحة (هجين)}};

  % SQL self-correction loop (max 3 retries)
  \draw[arrowA] (sql.north east)
    .. controls +(1.15,0.55) and +(1.15,-0.55) .. (sql.south east);
  \node[font=\tiny\bfseries, text=kaudark, align=center] at (8.95,-3.3)
    {\textarabic{حلقة تصحيح ذاتي}\\[-1pt]\textarabic{(حتى 3 محاولات)}};

  %==========================================================================
  % ANALYST (conditional call) and REPORTER
  %==========================================================================
  \node[agentbox] (analyst) at (6.4,-6.1)
    {\textarabic{وكيل التحليل \EN{(Analyst)}}\\[1pt]
     {\scriptsize\textarabic{إحصاء + كشف القيم الشاذة}}};
  \node[agentbox, minimum width=4.6cm] (rep) at (2.2,-8.35)
    {\textarabic{وكيل التقارير \EN{(Reporter)}}\\[1pt]
     {\scriptsize\textarabic{تنسيق \EN{RTL} + تصدير \EN{PDF/Excel/CSV}}}};

  \draw[arrowA] (sql.south) -- (analyst.north)
    node[lbls, pos=0.55, right=1pt] {\textarabic{بيانات خام}};
  \draw[arrowA] (analyst.south) -- (rep.north east)
    node[lbls, pos=0.55, right=2pt] {\textarabic{رؤى ومؤشرات}};

  % Reporter -> Supervisor final-result trunk (right side)
  \draw[arrowA] (rep.east) -- (10.2,-8.35) -- (10.2,-0.5) -- (sup.east);
  \node[font=\tiny\bfseries, text=kaudark!85, rotate=90] at (10.48,-4.4)
    {\textarabic{النتيجة النهائية}};

  %==========================================================================
  % Shared AgentState object + state bus (left side)
  %==========================================================================
  \node[statebox] (state) at (-5.3,-0.5)
    {\textarabic{كائن الحالة المشتركة \EN{AgentState}}\\[1pt]
     {\scriptsize\textarabic{يُقرأ ويُحدَّث في كل عقدة}}};
  \draw[statbus, <->, >=stealth] (state.east) -- (sup.west)
    node[lbls, midway, above=1pt] {\textarabic{قراءة / تحديث}};

  % Vertical state bus + horizontal stubs from each agent
  \draw[statbus] (-5.3,-1.15) -- (-5.3,-8.35);
  \draw[statbus] (rag.west) -- (-5.3,-3.3);
  \draw[statbus] (sql.south west) -- (4.0,-4.7) -- (-5.3,-4.7);
  \draw[statbus] (analyst.west) -- (-5.3,-6.1);
  \draw[statbus] (rep.west) -- (-5.3,-8.35);
  \node[font=\tiny\bfseries, text=kauprimary!80!black, rotate=90] at (-5.62,-5.2)
    {\textarabic{ناقل الحالة المشتركة \EN{(AgentState Bus)}}};

  %==========================================================================
  % Data stores: PostgreSQL (right) + Qdrant (left)
  %==========================================================================
  \node[datastore, minimum width=1.5cm, minimum height=1.0cm] (pg) at (9.2,-5.6)
    {\EN{PostgreSQL}};
  \draw[arrowL] (sql.south east) -- (pg.north);

  \node[datastore, minimum width=1.5cm, minimum height=0.95cm] (qd) at (-3.0,-5.4)
    {\EN{Qdrant}};
  \draw[arrowL] (rag.south west) -- (qd.north);

  %==========================================================================
  % Final outputs: answer pill (right) + audit box (left)
  %==========================================================================
  \node[outpill] (out) at (3.6,-10.05)
    {\textarabic{الإجابة النهائية للمستخدم}};
  \draw[arrowS] (rep.south) -- (out.north);

  \node[auditbox] (audit) at (-2.4,-10.05)
    {\textarabic{سجل التدقيق ومسار الوكلاء}\\[-1pt]{\tiny\EN{Audit Log + Agent Trace}}};
  \draw[arrowL] (rep.south west) -- (audit.north east)
    node[lbls, pos=0.45, left=2pt] {\textarabic{تسجيل}};

  %==========================================================================
  % Legend (bottom-right)
  %==========================================================================
  \node[draw=dividerdim!80, fill=white, rounded corners=3pt, inner sep=5pt,
        anchor=north east, font=\tiny, text=kaudark] at (10.3,-9.5)
  {%
    \begin{tabular}{@{}c@{\hspace{5pt}}r@{}}
      \tikz{\draw[arrowA] (0,0.08) -- (0.42,0.08);} & \textarabic{توجيه المشرف (حواف مشروطة)} \\[2pt]
      \tikz{\draw[arrowR] (0,0.08) -- (0.42,0.08);} & \textarabic{تدفق بيانات أو معامل بين الوكلاء} \\[2pt]
      \tikz{\draw[statbus] (0,0.08) -- (0.42,0.08);} & \textarabic{قراءة / تحديث \EN{AgentState}} \\[2pt]
      \tikz{\draw[arrowL] (0,0.08) -- (0.42,0.08);} & \textarabic{وصول إلى مخازن البيانات والتسجيل} \\
    \end{tabular}%
  };

\end{tikzpicture}%
}
}
  \caption[مخطط آلة الحالة وتدفق العقد المشروطة لمنظومة الوكلاء عبر LangGraph]{مخطط آلة الحالة وتدفق العقد المشروطة لمنظومة وكلاء نظام «بيان» عبر \EN{LangGraph} \EN{(Multi-Agent State \& Orchestration Graph)}: يوضح العقد الخمس، وكائن الحالة المشتركة \EN{AgentState} وناقله، والحواف المشروطة، وحلقة التصحيح الذاتي في وكيل \EN{SQL}، والمسار الهجين لتمرير المعامل من وكيل \EN{RAG}، والارتباط بطبقات البيانات وسجل التدقيق.}
  \label{fig:agent-state-graph}
\end{figure}

\subsection{نمذجة التسلسل التفاعلي الشامل للوكلاء}
\label{subsec:agent-sequence-diagram}

يكمل نمذجةَ المنظومة تتبعُ التسلسل الزمني لتنفيذ استعلام مؤسسي متكامل يمر بجميع الوكلاء، وهو ما يجسده مخطط التتابع \EN{(UML Sequence Diagram)} في الشكل \ref{fig:agent-sequence}. يوضح المخطط خطوط الحياة الثمانية مرتبة من اليمين إلى اليسار وفق اتجاه القراءة العربية: مستخدم المؤسسة، فواجهة المحادثة \EN{(Next.js)}، فالوكيل المشرف، فوكيل الاسترجاع \EN{RAG}، فوكيل قواعد البيانات عبر \EN{Smart Views}، فوكيل التحليل، فوكيل التقارير، وأخيراً وحدة سجل التدقيق وتتبع المسار.

ويعرض المخطط مثالاً تطبيقياً للنمط الهجين عبر سبع عشرة رسالة مرقمة موزعة على خمس مراحل: تبدأ مرحلة الاستيعاب والتحقق بتلقي الواجهة للسؤال العربي المركب وإرساله إلى المشرف وتسجيل بدء الجلسة، ثم تحليل النية وفحص الصلاحية ذاتياً؛ وتستخرج مرحلة استرجاع المعيار النسبة المسموحة من اللائحة عبر وكيل \EN{RAG} وبحثه المتجهي في \EN{Qdrant} وإعادتها موثقة المصدر؛ ثم تبني مرحلة استخراج البيانات الاستعلام الشرطي وتولّده وتصحيحه ذاتياً وتنفذه عبر وكيل \EN{SQL}؛ وتحسب مرحلة التحليل والتوليد المؤشرات والقيم الشاذة عبر وكيل التحليل ثم تنسق التقرير العربي عبر وكيل التقارير؛ وتختم مرحلة التسليم والتوثيق بتسجيل مسار الوكلاء كاملاً في سجل التدقيق وبث الإجابة النهائية إلى المستخدم عبر البث المتدفق \EN{(SSE)}. ويحقق هذا المخطط المهمة المجدولة في الخطة الزمنية للمشروع بالفصل الأول بعنوان «مخطط التسلسل الشامل»، ويوفر مرجعاً تنفيذياً دقيقاً لبناء الواجهات البرمجية بين المكونات في مرحلة التصميم.

\begin{figure}[H]
  \centering
  \resizebox{\textwidth}{!}{%==============================================================================
% figures/tikz/agent-sequence.tex
% Multi-Agent Interaction Sequence Diagram — new for Section 3.5 (2026-10).
%
% UML sequence diagram for one integrated hybrid query:
%   "أظهر الطلاب المتجاوزين لنسبة الغياب المسموحة وفق لائحة الجامعة"
% Lifelines ordered right-to-left (RTL reading): Tenant User -> Next.js UI ->
% Supervisor -> RAG Agent -> SQL Agent (Smart Views) -> Analyst -> Reporter ->
% Audit & Trace. 17 numbered messages, 5 phase bands, activation bars on the
% supervisor / UI / specialist agents. Solid arrows = calls, dashed =
% returns, self-loops = internal processing (incl. SQL self-correction).
% All Arabic labels wrapped in \textarabic{} (proven RTL fix).
%==============================================================================
\resizebox{\textwidth}{!}{%
\begin{tikzpicture}[
  >=stealth,
  font=\small,
  x=1cm, y=1cm,
  head/.style={
    rectangle, rounded corners=4pt, draw=kauprimary, fill=kauprimary!10,
    text=kaudark, align=center, inner sep=4pt,
    font=\scriptsize\bfseries, minimum width=2.25cm, minimum height=0.82cm
  },
  call/.style={->, >=stealth, line width=0.75pt, draw=kaudark!80},
  ret/.style={->, >=stealth, line width=0.6pt, draw=kausecondary!90!black, dashed},
  selfm/.style={->, >=stealth, line width=0.6pt, draw=kaudark!75},
  lifeline/.style={draw=kaudark!45, dashed, line width=0.5pt},
  phase/.style={font=\tiny\bfseries, text=kauprimary!85!black},
  msgl/.style={font=\tiny, text=kaudark, fill=white, inner sep=0.8pt},
  actbar/.style={fill=kausecondary!30, draw=kausecondary!80!black, line width=0.4pt},
]

  %==========================================================================
  % Lifeline x-positions (right -> left = RTL reading order)
  %==========================================================================
  \def\ux{8.8}     % Tenant User
  \def\uix{6.45}   % Next.js chat UI
  \def\sux{4.1}    % Supervisor Agent
  \def\ragx{1.75}  % RAG Agent
  \def\sqlx{-0.6}  % SQL Agent (Smart Views)
  \def\anx{-2.95}  % Analyst Agent
  \def\repx{-5.3}  % Reporter Agent
  \def\aux{-7.65}  % Audit & Trace

  % ---- Phase bands (alternating tints, drawn first) ----
  \fill[kauprimary!4]  (10.6,-1.05)  rectangle (-8.7,-4.30);
  \fill[kausecondary!5] (10.6,-4.30)  rectangle (-8.7,-7.20);
  \fill[kauprimary!4]  (10.6,-7.20)  rectangle (-8.7,-9.80);
  \fill[kausecondary!5] (10.6,-9.80)  rectangle (-8.7,-12.90);
  \fill[kauprimary!4]  (10.6,-12.90) rectangle (-8.7,-15.35);

  % ---- Phase separators + labels (right margin = first read) ----
  \foreach \yy in {-4.30,-7.20,-9.80,-12.90}{%
    \draw[draw=dividersilver, dotted, line width=0.5pt] (10.6,\yy) -- (-8.7,\yy);
  }
  \node[phase] at (9.75,-1.35)  {\textarabic{المرحلة 1: الاستيعاب والتحقق}};
  \node[phase] at (9.75,-4.60)  {\textarabic{المرحلة 2: استرجاع المعيار من اللائحة}};
  \node[phase] at (9.75,-7.50)  {\textarabic{المرحلة 3: استخراج البيانات}};
  \node[phase] at (9.75,-10.10) {\textarabic{المرحلة 4: التحليل وتوليد التقرير}};
  \node[phase] at (9.75,-13.20) {\textarabic{المرحلة 5: التسليم والتوثيق}};

  % ---- Lifelines ----
  \foreach \xx in {8.8,6.45,4.1,1.75,-0.6,-2.95,-5.3,-7.65}{%
    \draw[lifeline] (\xx,-0.82) -- (\xx,-15.35);
  }

  % ---- Header boxes ----
  \node[head] at (\ux,0)    {\textarabic{مستخدم المؤسسة}\\[-1pt]{\tiny\EN{(Tenant User)}}};
  \node[head] at (\uix,0)   {\textarabic{واجهة المحادثة}\\[-1pt]{\tiny\EN{Next.js}}};
  \node[head, draw=kauprimary, fill=kauprimary!16, line width=0.9pt] at (\sux,0)
    {\textarabic{الوكيل المشرف}\\[-1pt]{\tiny\EN{(Supervisor)}}};
  \node[head] at (\ragx,0)  {\textarabic{وكيل الاسترجاع}\\[-1pt]{\tiny\EN{(RAG Agent)}}};
  \node[head] at (\sqlx,0)  {\textarabic{وكيل قواعد البيانات}\\[-1pt]{\tiny\EN{SQL + Smart Views}}};
  \node[head] at (\anx,0)   {\textarabic{وكيل التحليل}\\[-1pt]{\tiny\EN{(Analyst)}}};
  \node[head] at (\repx,0)  {\textarabic{وكيل التقارير}\\[-1pt]{\tiny\EN{(Reporter)}}};
  \node[head, draw=kaudark!70, fill=kaudark!6] at (\aux,0)
    {\textarabic{التدقيق والتتبع}\\[-1pt]{\tiny\EN{Audit \& Trace}}};

  % ---- Activation bars ----
  \fill[actbar] (\uix-0.07,-1.45) rectangle (\uix+0.07,-14.55);  % UI (streaming)
  \fill[actbar] (\sux-0.09,-2.20) rectangle (\sux+0.09,-14.25);  % Supervisor
  \fill[actbar] (\ragx-0.07,-4.70) rectangle (\ragx+0.07,-6.30); % RAG
  \fill[actbar] (\sqlx-0.07,-7.55) rectangle (\sqlx+0.07,-9.15); % SQL
  \fill[actbar] (\anx-0.07,-9.85) rectangle (\anx+0.07,-10.70);  % Analyst
  \fill[actbar] (\repx-0.07,-11.40) rectangle (\repx+0.07,-12.25);% Reporter
  \fill[actbar] (\aux-0.07,-2.92) rectangle (\aux+0.07,-3.10);   % Audit tick 1
  \fill[actbar] (\aux-0.07,-12.92) rectangle (\aux+0.07,-13.10); % Audit tick 2

  %==========================================================================
  % Messages
  %==========================================================================
  % Phase 1: intake & verification
  \draw[call] (\ux,-1.50) -- (\uix,-1.50)
    node[msgl, midway, above] {\textarabic{1: استعلام عربي مركّب (بيانات + لائحة)}};
  \draw[call] (\uix,-2.30) -- (\sux,-2.30)
    node[msgl, midway, above] {\textarabic{2: إرسال الطلب \EN{(POST /api/query)}}};
  \draw[call] (\sux,-3.00) -- (\aux,-3.00)
    node[msgl, midway, above] {\textarabic{3: تسجيل بدء الجلسة \EN{(session id)}}};
  \draw[selfm] (\sux,-3.55) .. controls +(0.75,0.0) and +(0.75,-0.52) .. (\sux,-4.07);
  \node[msgl, anchor=west] at (\sux+0.85,-3.81) {\textarabic{4: تحليل النية + فحص \EN{RBAC}}};

  % Phase 2: retrieve the regulatory threshold
  \draw[call] (\sux,-4.80) -- (\ragx,-4.80)
    node[msgl, midway, above] {\textarabic{5: استخراج نسبة الغياب المسموحة من اللائحة}};
  \draw[selfm] (\ragx,-5.35) .. controls +(0.75,0.0) and +(0.75,-0.5) .. (\ragx,-5.85);
  \node[msgl, anchor=west] at (\ragx+0.85,-5.60) {\textarabic{6: بحث متجهي في \EN{Qdrant}}};
  \draw[ret] (\ragx,-6.45) -- (\sux,-6.45)
    node[msgl, midway, above] {\textarabic{7: النسبة = 25\% مع توثيق المصدر}};

  % Phase 3: structured data extraction
  \draw[call] (\sux,-7.65) -- (\sqlx,-7.65)
    node[msgl, midway, above] {\textarabic{8: استعلام: الطلاب بغياب يفوق 25\%}};
  \draw[selfm] (\sqlx,-8.20) .. controls +(0.75,0.0) and +(0.75,-0.5) .. (\sqlx,-8.70);
  \node[msgl, anchor=west] at (\sqlx+0.85,-8.45) {\textarabic{9: توليد \EN{SQL} + تصحيح ذاتي}};
  \draw[ret] (\sqlx,-9.30) -- (\sux,-9.30)
    node[msgl, midway, above] {\textarabic{10: بيانات خام \EN{(DataFrame)}}};

  % Phase 4: analysis & report generation
  \draw[call] (\sux,-10.00) -- (\anx,-10.00)
    node[msgl, midway, above] {\textarabic{11: تحليل إحصائي وكشف شواذ}};
  \draw[ret] (\anx,-10.62) -- (\sux,-10.62)
    node[msgl, midway, above] {\textarabic{12: مؤشرات + تنبيهات شذوذ}};
  \draw[call] (\sux,-11.50) -- (\repx,-11.50)
    node[msgl, midway, above] {\textarabic{13: توليد التقرير المنسق \EN{(RTL)}}};
  \draw[ret] (\repx,-12.12) -- (\sux,-12.12)
    node[msgl, midway, above] {\textarabic{14: تقرير نهائي + تصدير \EN{PDF/Excel}}};

  % Phase 5: delivery & auditing
  \draw[call] (\sux,-13.00) -- (\aux,-13.00)
    node[msgl, midway, above] {\textarabic{15: تسجيل مسار الوكلاء \EN{(Agent Trace)}}};
  \draw[call] (\sux,-13.80) -- (\uix,-13.80)
    node[msgl, midway, above] {\textarabic{16: بث الإجابة \EN{(SSE Stream)}}};
  \draw[call] (\uix,-14.50) -- (\ux,-14.50)
    node[msgl, midway, above] {\textarabic{17: عرض التقرير والرسوم والتصدير}};

\end{tikzpicture}%
}
}
  \caption[مخطط التسلسل التفاعلي الشامل لمنظومة وكلاء نظام «بيان»]{مخطط التسلسل التفاعلي الشامل \EN{(UML Sequence Diagram)} لتنفيذ استعلام مؤسسي متكامل عبر منظومة وكلاء نظام «بيان»: خطوط الحياة الثمانية من اليمين إلى اليسار، والرسائل السبع عشرة الموزعة على خمس مراحل (الاستيعاب والتحقق، استرجاع المعيار من اللائحة، استخراج البيانات، التحليل وتوليد التقرير، التسليم والتوثيق).}
  \label{fig:agent-sequence}
\end{figure}

